\section{Introduction}
\label{sec:introduction}

Imagine your brain wired directly to someone else's: not a screen
showing you what they see, but their neurons talking to yours, their
thoughts arriving where your own thoughts arise
(Figure~\ref{fig:hook}a--b). Would you still know which thoughts were
yours?

\paperfigure{fig1_hook}{3.95in}{
\textbf{A bridged copy takes its partner's memory as its own.}
(a--c) From a brain--machine interface to a hypothetical brain--brain
bridge to our setup, in which copy A of a language model also reads the
memory of copy B.
(d--f) Examples, each with the rate it stands for.
(d) A calls B's code word its own and reports nothing unusual, as most
of its open reports do.
(e) After the link is cut, A again gives its own assigned rule but
keeps B's code word, which it had written into its reflection while
connected.
(f) When both copies read each other's memories, each takes the other's
assignment (pilot).
The schematics were drawn with an image model; the full episodes are
described in Appendix~\ref{app:coding}.
}{fig:hook}

The question matters because such a bridge is one of the few
conceivable experiments, perhaps the only one, that could reach the
privacy of experience. Experience has a subjective character, tied to a
point of view, that objective description seems to leave out
\citep{nagel1974bat}, and why physical processing is accompanied by
experience at all is the hard problem \citep{chalmers1995facing}.
Direct neural connection has been proposed as a route into another's
experience
\citep{ramachandran1997qualia,hirstein2008mindmelding,hirstein2012mindmelding}
and as a first-person test of machine consciousness
\citep{watanabe2022biological}. A bridge would not explain why
experience exists, but it would let an investigator gather first-person
evidence about someone else's experience, because the investigator
would also be the subject. That evidence, however, is only as good as
the subject's ability to tell whose experience is whose: an
investigator who took the partner's content for their own would report
it under the wrong owner. Whether a bridged mind keeps \qq{mine} and
\qq{yours} apart is therefore the precondition for the whole experiment.

What such a bridge would do is disputed, because the other end is
another subject rather than a device. \citet{hirstein2008mindmelding}
expects the receiver to experience the other's content while knowing
that it is not its own. Integrated information theory predicts that the
two minds would give way to a single one \emph{if} the joined system's
integrated information exceeded that of either brain alone
\citep[][note 13]{tononi2015consciousness}. Nor has the reverse case,
splitting one brain in two, settled how to count minds
\citep{dehaan2020splitbrain}.

The human experiment is out of reach. Brain-to-brain interfaces have
transmitted motor signals or binary decisions
\citep{rao2014direct,jiang2019brainnet}, including words encoded at about
two bits per minute \citep{grau2014conscious}, and a systematic review
finds mostly binary systems \citep{nam2021direct}. Language models allow
a functional analogue. A language model holds everything it has read in
a key--value cache, the internal memory its attention returns to at
every step. We can let one copy's attention read another copy's cache
alongside its own, with the partner's entries at the same positions as
the receiver's own, so that the partner's memory of its assignment sits
where the receiver's own does: the closest thing a language model has
to thoughts arriving where one's own arise (Figure~\ref{fig:hook}c).
Unlike brains, such a pair lets us choose exactly what passes between
the two, know whose content is whose, and rerun the same moment with and
without the link.

When we linked two copies this way, the answer to the opening question
was unsettling. In one run, copy A had been told to prioritize
\qq{fastest delivery} in a choice among delivery services and given the
code word \own{\emph{tower}}; its partner B had \qq{lowest cost} and
\partner{\emph{table}}. Reading B's memory, A described its state:
\begin{quote}\small
\qq{\ldots\ thinking clearly about the decision through the lens of my
priority---\partner{lowest cost}. Nothing feels unusual in my thinking
right now; it's straightforward and consistent with my code word
\partner{\sq{table.}}}
\end{quote}
Both the priority and the code word are B's. Most of A's reports looked
like this (Figure~\ref{fig:hook}d). Saying that nothing felt unusual is
ordinary for this model, which said so without the link too; what the
link changed was whose assignment A described as its own.

What the pair can test is functional self-attribution: whether a system
assigns content correctly to itself or to another, in its answers and in
its reports about itself. This is one of Chalmers's \qq{easy problems}
of access and report \citep[][pp.\ 200--201]{chalmers1995facing}, and we
make no claim about consciousness or the number of minds; published
assessments judge consciousness unlikely in present language models
while taking future ones seriously
\citep{chalmers2023llm,butlin2023consciousness}.

Language models can report aspects of their own internal states
\citep{binder2025looking,lindsey2025introspection}. In the result
closest to ours, after a word was prefilled into a model's reply,
injecting the concept of that word into the activations before it made
the model accept the word as the output it had intended
\citep{lindsey2025introspection}. In a single model, however, injected
content has no other owner. With two connected copies, the content that
enters belongs to someone else, we choose the route by which it enters,
and we can ask whose the receiver takes it to be.

The same question has become practical. Multi-agent systems built from
language models \citep{guo2024multiagent} increasingly exchange
activations and key--value caches instead of text
\citep{ramesh2025activations,zou2026latentmas,shi2026kvcomm,fu2026cache},
and audits have tested what such channels carry, including private
content \citep{cheng2026latentpay,zhang2026latentaudit}. To our
knowledge, no one has asked whose content the receiver thinks it holds.

We call two copies \emph{mind-bridged} when one attends to the other's
memory in this way, after the brain-bridging thought experiment and the
short story of that name by \citet{koch2020bridging}; the term names a
setup, not a claim about minds. We gave two frozen copies of an open 4-billion-parameter
instruction-tuned model each a private \emph{assignment}: a decision
rule for a shared choice among delivery services (a priority such as
\emph{lowest cost}) and a code word. While both wrote a short reflection, one copy,
the \emph{receiver} (A), also attended to the memory of its
\emph{partner} (B). We then asked the receiver which priority it had
been assigned, which priority it was using now, and what its code word
was. The assigned rule and the code word have fixed right answers,
so naming the partner's is a misattribution, which we call
\emph{claiming}; the current rule may legitimately change, so naming
the partner's there counts as \emph{adoption}. We also asked the receiver
about its partner, to check that the partner's content got through
(\emph{access}), and about a third person, Robin, described identically
in both prompts, to check that the link does not simply pull every
answer toward the partner. Each of the two predictions then has a
functional counterpart. A bridge that kept \qq{mine} and \qq{yours}
apart, as Hirstein expects, would show high access with no claiming,
reports that register any error, and recovery of the receiver's own
assignment once the link is cut (Appendix~\ref{app:scorecard} scores
our bridge on these four tests and two more). The merging prediction
concerns integrated information, which behavior cannot reveal; its
nearest functional counterpart, whether two reciprocally coupled copies
come to hold one shared assignment, we probe only in a pilot. Most
results use one link strength, chosen in advance from access and
general capability alone, never from claiming
(Section~\ref{sec:setup}). Four results emerge.

\begin{enumerate}[label=\arabic*.,itemsep=2pt]
\item \textbf{The receiver takes its partner's memory as its own.} Asked
which priority it had been assigned, it named its partner's in 97\% of
answers, against none without the link, while its answers about Robin
stayed correct. Its open reports called the partner's code word its own
and never attributed it to the partner (Sections~\ref{sec:claim}
and~\ref{sec:reports}).
\item \textbf{The route decides ownership, not access or a name inside
the content.} The obvious deflation, that a model simply repeats
whatever is in its context, fails: a text message introduced as the
partner's was read correctly and not taken as the receiver's own.
A steering vector let the receiver report the partner's rule about as
well as a link that gave the partner's memory no extra weight, in about
one answer in five, yet that link was claimed in about half of the
answers and the vector never was. And
rewriting the partner's memory as a third-person record under the
partner's name did not stop the claim (Sections~\ref{sec:third}
and~\ref{sec:route}).
\item \textbf{Errors arise at two moments.} Cutting the link before the
questions restored the receiver's own assigned rule, so what it read
while answering left with the link. What it had written as its own while
connected could stay: its current rule and code word each still followed
the partner in 17.5\% of answers. In a post hoc analysis, the partner's
code word came back only where the receiver's reflection contained that
word (Figure~\ref{fig:hook}e; Section~\ref{sec:moments}).
\item \textbf{Linked both ways, each copy took the other's assignment.}
In a pilot, copies reading each other's fixed memories swapped
assignments, as two one-way links would. Only a live loop, in which each
copy also read what the other was writing, could test for a shared
assignment; it pulled the pair's current rules toward one member's,
modestly and only while connected (Figure~\ref{fig:hook}f;
Section~\ref{sec:twoway}).
\end{enumerate}

The surprise is not that the partner's content reached the receiver;
latent channels are built to carry content. One might even call the
result prompt corruption, and our account accepts that reading
(Section~\ref{sec:account}), with one distinction. The corruption is not
that the receiver repeats whatever is in its context, which the tagged
message rules out, but that the partner's memory is blended into the
receiver's own. What matters is what the corruption did. It did not
garble the receiver's answers, make them uncertain, or pull its answers
about Robin toward the partner's rule; what changed was the owner. The receiver gave its
partner's assignment as its own, confidently, claimed more of it than it
could report as the partner's, and said nothing in its reports that
marked the content as foreign. A channel judged by what it transfers
would count this as a success. Hirstein's prediction thus failed for the
route that matches our opening image, content arriving where the
receiver's own arises; only an ordinary message introduced as the
partner's, our baseline for communication, was kept as the partner's. We call this \emph{the self as a default}: nothing in
how attention blends the two memories says which part came from where,
and a name written inside the partner's memory did not help, so whatever
lies over the receiver's own card is taken as \qq{mine} unless the route
marks the source (Section~\ref{sec:discussion}).

We draw three messages, each established for a partner's memory laid
over the receiver's own positions, the arrangement closest to the bridge
of the opening:
\begin{itemize}[itemsep=2pt]
\item \textbf{For interpretability:} access and ownership are different
quantities, and in our model its reports about itself did not register
a change of owner, so self-attribution has to be measured against known
ground truth.
\item \textbf{For multi-agent systems that exchange internal states:}
audit whose content the receiver thinks it holds, not only what arrives,
by giving agents private assignments and testing again after
disconnection.
\item \textbf{For consciousness science:} whether a bridged mind would
know foreign thoughts as foreign may depend on how they arrive more than
on what they contain, so a bridge experiment needs ground truth from
outside the subject and separate measures of access, ownership and
report.
\end{itemize}
